\section{传输、乘积与一维分布}\label{sec:02-03}

\subsection{可测映射、推前测度与分布}\label{subsec:2-3-1}
\index{推前测度}\index{分布}

一个测度未必直接写在我们关心的空间上。样本经过统计量变成一个数，Euclidean 点经过坐标投影
变成一个坐标，几何区域经过变换被搬到另一处。目标集合 $B$ 的质量只能通过源中的逆像
$T^{-1}(B)$ 计算；这正是推前的定义。它保留由 $T$ 能观察到的信息，也会不可逆地丢掉同一纤维内的
差别。

\begin{definition}[可测映射]\label{def:2-3-1-1}
映射 $T:(X,\mathcal A)\to(Y,\mathcal B)$ 称为可测映射，如果
$T^{-1}(B)\in\mathcal A$ 对每个 $B\in\mathcal B$ 成立。若 $T$ 为双射且 $T,T^{-1}$ 都可测，
称为可测同构。
\end{definition}

\begin{definition}[推前测度]\label{def:2-3-1-2}
给定 $X$ 上测度 $\mu$ 和可测映射 $T:X\to Y$，定义
\[
 T_*\mu(B)=\mu(T^{-1}(B)),\qquad B\in\mathcal B.
\]
若 $\mu(X)=1$，称 $(X,\mathcal A,\mu)$ 为概率空间，其中的可测映射暂称随机变量，并把
$T_*\mu$ 记作 $\Law(T)$。若可测同构还满足 $T_*\mu=\nu$，称其为测度空间同构。
\end{definition}

逆像保持空集、补集和可数并，因此 $T_*\mu$ 的可数可加性由 $\mu$ 的可数可加性逐项得到。

\begin{proposition}[函子性]\label{proposition:2-3-1-1}
若复合有意义，则
\[
 (S\circ T)_*\mu=S_*(T_*\mu),\qquad (\id_X)_*\mu=\mu.
\]
\end{proposition}

\begin{proof}
对目标可测集 $C$，
\[
 (S\circ T)_*\mu(C)
 =\mu(T^{-1}(S^{-1}C))
 =T_*\mu(S^{-1}C)=S_*(T_*\mu)(C).
\]
恒等映射的逆像不改变集合，第二式随之成立。
\end{proof}

\begin{proposition}[有限简单测试公式]\label{proposition:2-3-1-1a}
若 $B_1,\dots,B_r\in\mathcal B$ 两两不交且 $a_k\ge0$，则
\[
 \sum_{k=1}^ra_k(T_*\mu)(B_k)
 =\sum_{k=1}^ra_k\mu(T^{-1}B_k).
\]
\end{proposition}

\begin{proof}
每一项由推前定义相等；逆像还保持两两不交，所以两边确实是同一个有限分割上的质量和。
\end{proof}

\begin{proposition}[连续推前与支撑]\label{proposition:2-3-1-1b}
设 $X,Y$ 为拓扑空间，$\mu$ 为 $X$ 上的 Borel 测度，且
$\mu(X\setminus\supp\mu)=0$。若 $T:X\to Y$ 连续，则
\[
 \supp(T_*\mu)=\overline{T(\supp\mu)}.
\]
\end{proposition}

\begin{proof}
若 $x\in\supp\mu$，含 $T(x)$ 的每个开集 $V$ 有开逆像 $T^{-1}V$ 含 $x$，故
$T_*\mu(V)=\mu(T^{-1}V)>0$。所以 $T(\supp\mu)\subset\supp(T_*\mu)$；右侧支撑闭，得到一个
包含关系。

若 $y\notin\overline{T(\supp\mu)}$，取开邻域
$V=Y\setminus\overline{T(\supp\mu)}$。它的逆像与 $\supp\mu$ 不交，因此
$T_*\mu(V)\le\mu(X\setminus\supp\mu)=0$，所以 $y$ 不在推前支撑中。
\end{proof}

\begin{theorem}[推前积分公式]\label{theorem:2-3-1-2}
对非负可测函数 $f:Y\to[0,\infty]$，
\[
 \int_Y f\,\dd(T_*\mu)=\int_Xf\circ T\,\dd\mu.
\]
\end{theorem}

非负积分将在下一章由简单函数下逼近定义，届时
\cref{proposition:3-1-1-5} 将证明上式；可积实值和复值版本分别由正负部与实虚部得到。

\begin{proposition}[连续推前保持 Radon]\label{proposition:2-3-1-3}
设 $X,Y$ 为局部紧 Hausdorff 空间，$\mu$ 为有限 Radon 测度，$T:X\to Y$ 连续，则
$T_*\mu$ 为 Radon。若 $\mu$ 不必有限而 $T$ 为固有映射（proper map），即每个紧集的逆像紧，
则结论仍成立。
\end{proposition}

\begin{proof}
先记录有限 Radon 测度的一个推论：它对每个 Borel 集都可由紧集从内逼近。事实上，若
$E\subset X$ 为 Borel 集且 $\varepsilon>0$，由开集 $X$ 的内正则性可取紧集 $K_0\subset X$，使
$\mu(X\setminus K_0)<\varepsilon/2$；又由 $E^c$ 的外正则性可取开集 $O\supset E^c$，使
$\mu(O)<\mu(E^c)+\varepsilon/2$。由于 $\mu$ 有限，
\[
 K:=K_0\setminus O\subset E,\qquad
 \mu(E\setminus K)
 \le \mu(X\setminus K_0)+\mu(O\setminus E^c)<\varepsilon.
\]
这里 $K$ 是紧集。

现在先设 $\mu$ 有限，并记 $\nu=T_*\mu$。若 $B\subset Y$ 为 Borel 集，对 $T^{-1}B$ 应用上面的
紧内逼近，得到紧集 $K\subset T^{-1}B$，使
$\mu(T^{-1}B\setminus K)<\varepsilon$；其像 $T(K)\subset B$ 紧且
\[
 \nu(T(K))=\mu(T^{-1}(T(K)))\ge\mu(K).
\]
从而
\[
 0\le \nu(B)-\nu(T(K))
 \le \mu(T^{-1}B)-\mu(K)
 =\mu(T^{-1}B\setminus K)<\varepsilon.
\]
这就得到 $\nu$ 对 Borel 集的紧内逼近。再对 $B^c$ 作
紧内逼近：若紧集 $L\subset B^c$ 满足 $\nu(B^c\setminus L)<\varepsilon$，则开集
$Y\setminus L$ 包含 $B$，并且
\[
 \nu((Y\setminus L)\setminus B)=\nu(B^c\setminus L)<\varepsilon.
\]
这给出外正则性。最后，目标紧集的质量至多 $\nu(Y)=\mu(X)<\infty$，故 $\nu$ 是 Radon 测度。

再设 $T$ 固有。目标紧集 $L$ 的逆像紧，故
$T_*\mu(L)=\mu(T^{-1}L)<\infty$。还需证明正则性。局部紧 Hausdorff 空间之间的固有连续映射是闭
映射：若 $F$ 闭、$y\notin T(F)$，取含 $y$ 的相对紧开邻域 $V$。集合
$F\cap T^{-1}(\overline V)$ 是紧集，故其像
$T(F\cap T^{-1}(\overline V))$ 是 $\overline V$ 中的闭集，而且不含 $y$。因此可在 $V$ 中取
$y$ 的开邻域 $W$，使它与这个像不交。若 $z\in W\cap T(F)$，则
$z\in\overline V$，其某个 $F$ 中原像会落在 $F\cap T^{-1}(\overline V)$，矛盾。所以
$W\cap T(F)=\varnothing$，从而 $T(F)$ 闭。

仍记 $\nu=T_*\mu$。若 $V\subset Y$ 开，则 $T^{-1}V$ 开。若
$\mu(T^{-1}V)<\infty$，给定 $\varepsilon>0$，源测度的开集内正则性给出紧集
$K\subset T^{-1}V$，使 $\mu(K)>\mu(T^{-1}V)-\varepsilon$；若
$\mu(T^{-1}V)=\infty$，则对每个 $M>0$ 可取这样的紧集使 $\mu(K)>M$。两种情形中
$T(K)\subset V$ 都是紧集，且 $\nu(T(K))\ge\mu(K)$。因此 $\nu$ 对开集内正则。

再证外正则。给定 Borel 集 $B\subset Y$；若 $\nu(B)=\infty$，外正则等式由单调性直接成立。
若 $\nu(B)<\infty$，源测度的外正则性给出开集 $O\supset T^{-1}B$，使
\[
 \mu(O)<\mu(T^{-1}B)+\varepsilon.
\]
闭映射性使
$V=Y\setminus T(X\setminus O)$ 为含 $B$ 的开集，且 $T^{-1}V\subset O$，所以
\[
 \nu(V)=\mu(T^{-1}V)\le\mu(O)<\nu(B)+\varepsilon.
\]
这证明了外正则性；目标紧集的有限性已在上面证明，故 $\nu$ 是 Radon 测度。
\end{proof}

\begin{proposition}[分布的集合层识别]\label{proposition:2-3-1-4}
若随机变量 $\xi,\eta$ 同分布，则
$\Prob(\xi\in B)=\Prob(\eta\in B)$ 对每个目标 Borel 集成立。反过来，若该等式在一个生成
Borel $\sigma$-代数的 $\pi$-系统上成立，则 $\xi,\eta$ 同分布。
\end{proposition}

\begin{proof}
第一句是推前测度相等的展开。反向中两个分布都是概率测度，故在生成 $\pi$-系统上相等后可用
\cref{theorem:2-2-1-2}。
\end{proof}

\begin{example}[经验测度]\label{ex:2-3-1-1}
给定 $x_1,\dots,x_n\in X$，在有限集合 $\{1,\dots,n\}$ 上取均匀概率，令 $T(k)=x_k$，则
\[
 T_*\Bigl(\frac1n\sum_{k=1}^n\delta_k\Bigr)
 =\frac1n\sum_{k=1}^n\delta_{x_k}=:L_n.
\]
对有限简单测试函数，两边都化为有限样本和；经验测度因此把有限数据变成测度对象。
\end{example}

\begin{example}[坐标分布]\label{ex:2-3-1-2}
若 $\mu$ 是 $\R^n$ 上测度，坐标映射 $x\mapsto x_j$ 的推前给出第 $j$ 坐标分布；线性函数
$x\mapsto\langle t,x\rangle$ 的推前记录沿方向 $t$ 的一维观测。例如在 \(\R^2\) 上取
\[
 \mu=\frac12\delta_{(0,1)}+\frac12\delta_{(2,-1)}.
\]
若 \(\pi_1,\pi_2\) 是两个坐标投影，则
\[
 (\pi_1)_*\mu=\frac12\delta_0+\frac12\delta_2,\qquad
 (\pi_2)_*\mu=\frac12\delta_1+\frac12\delta_{-1}.
\]
对 \(t=(1,1)\)，两点的内积分别为 \(1,1\)，所以
\((x\mapsto\langle t,x\rangle)_*\mu=\delta_1\)。\Cramer{}--Wold 方法将研究全部这类一维观测
是否足以识别高维分布。
\end{example}

\begin{example}[非单射丢失信息]\label{ex:2-3-1-3}
$T(x)=x^2$ 满足 $T_*\delta_1=T_*\delta_{-1}=\delta_1$。推前不能区分落在同一纤维中的两个点，
所以没有单射条件便不能反演。
\end{example}

\begin{example}[平方折叠的集合计算]\label{ex:2-3-1-4}
令 $\mu=\tfrac12m|_{[-1,1]}$，$T(x)=x^2$。若 $0\le a<b\le1$，则
\[
 T^{-1}((a,b])=[-\sqrt b,-\sqrt a)\cup(\sqrt a,\sqrt b],
\]
除端点约定外两段长度都为 $\sqrt b-\sqrt a$，所以
\[
 (T_*\mu)((a,b])=\sqrt b-\sqrt a.
\]
第~3~章的换元公式将由这个增量导出推前测度的密度。
\end{example}

\begin{figure}[htbp]
\centering
\begin{tikzpicture}[node distance=17mm and 20mm]
 \node[book strong node,minimum width=38mm] (x) {源空间 $(X,\mu)$\\纤维 $T^{-1}(y)$};
 \node[book strong node,right=of x,minimum width=38mm] (y) {目标空间 $(Y,T_*\mu)$\\集合 $B$};
 \draw[book accent arrow] (x)--node[above]{$T$}(y);
 \draw[book dashed arrow,bend right=22] (y.south)--node[below]{$B\mapsto T^{-1}(B)$}(x.south);
\end{tikzpicture}
\caption{推前沿 $T$ 移动质量，却以逆像计算目标集合。}
\label{fig:pushforward-fibres}
\end{figure}

\begin{remark}
可测映射只要求可测集的逆像可测；可测集的像未必可测。同分布也不表示两个随机变量在同一概率
空间上相等。无限 Radon 测度的推前若缺少固有性，目标紧集的逆像可能携带无穷质量，局部
有限性便会失败。
\end{remark}

对 \(x\mapsto\langle t,x\rangle\) 的推前取复指数积分，便得到特征函数；因此方向投影的一维分布
是 \Cramer{}--Wold 方法的基本数据。Markov 核把“每个起点对应一个目标测度”参数化，其确定性特例
正是 \(x\mapsto\delta_{T(x)}\)。这些概念将在概率论部分定义。

\begin{exercise}
从逆像恒等式证明推前函子性。
\end{exercise}
\begin{exercise}
计算仿射映射对有限原子测度和区间均匀概率的推前。
\end{exercise}
\begin{exercise}
分别写出支撑公式两个包含关系中使用的假设，并给出删除连续性后的反例。
\end{exercise}
\begin{exercise}
补全固有连续映射为闭映射的证明，并据此证明无限 Radon 推前命题。
\end{exercise}
\begin{exercise}
在同一概率空间上构造同分布但不几乎处处相等的两个随机变量。
\end{exercise}

单个映射只搬运一份质量。要同时描述两个观测及其联合事件，还需在矩形上规定质量，并证明这份有限数据
能够扩张到乘积 $\sigma$-代数。

\subsection{乘积 \texorpdfstring{$\sigma$}{sigma}-代数、截面与乘积测度}\label{subsec:2-3-2}
\index{乘积测度}\index{截面}

两个空间中的有限观测组合成矩形 $A\times B$。若规定矩形质量为 $\mu(A)\nu(B)$，存在性需要
把这个规则扩张到可数生成事件，唯一性需要证明矩形足以识别扩张。下面直接在有限简单函数的质量和
上验证预测度；第~3~章的 Tonelli 定理将以乘积测度的存在为基础。

\begin{definition}[乘积 $\sigma$-代数]\label{def:2-3-2-1}
$\mathcal A\otimes\mathcal B$ 是所有可测矩形 $A\times B$ 生成的 $\sigma$-代数。对任意指标族，
柱 $\sigma$-代数由只限制有限多个坐标的坐标逆像生成。
\end{definition}

\begin{definition}[截面]\label{def:2-3-2-2}
对 $E\subset X\times Y$，定义
\[
 E_x=\{y:(x,y)\in E\},\qquad E^y=\{x:(x,y)\in E\}.
\]
对函数 $f$，记 $f_x(y)=f(x,y)$。
\end{definition}

\begin{proposition}[截面可测]\label{proposition:2-3-2-1}
若 $E\in\mathcal A\otimes\mathcal B$，则每个 $E_x\in\mathcal B$、每个
$E^y\in\mathcal A$。
\end{proposition}

\begin{proof}
固定 $x$，令 $\mathcal C_x=\{E:E_x\in\mathcal B\}$。矩形的截面是 $B$ 或空集；截面运算保持
补集和可数并，所以 $\mathcal C_x$ 是含所有矩形的 $\sigma$-代数，遂包含
$\mathcal A\otimes\mathcal B$。固定 $y$ 的证明相同。
\end{proof}

\begin{example}[逐截面可测仍不够]\label{ex:2-3-2-4}
取 $A\subset X$ 为 $\mathcal A$-不可测集，取 $B\in\mathcal B$ 且
$0<\nu(B)<\infty$，令 $E=A\times B$。每个竖截面是 $B$ 或空集，因而可测；但若
$y\in B$，水平截面 $E^y=A$ 不可测，所以上一命题的逆否命题给出
$E\notin\mathcal A\otimes\mathcal B$。而且
$x\mapsto\nu(E_x)=\nu(B)\1_A(x)$ 不可测。逐条截面可测与联合可测是两件事。
\end{example}

\begin{definition}[边缘测度]\label{def:2-3-2-4}
若 $\pi$ 是 $(X\times Y,\mathcal A\otimes\mathcal B)$ 上的测度，坐标投影可测；推前
$(\pi_X)_*\pi$ 与 $(\pi_Y)_*\pi$ 称为 $\pi$ 的两个边缘测度。固定边缘寻找联合测度称为耦合
问题。
\end{definition}

\begin{definition}[乘积测度]\label{def:2-3-2-3}
对 $\sigma$-有限测度 $\mu,\nu$，满足
\[
 (\mu\otimes\nu)(A\times B)=\mu(A)\nu(B),
 \qquad A\in\mathcal A, B\in\mathcal B,
\]
的测度称为乘积测度；约定 $0\cdot\infty=0$。
\end{definition}

$\sigma$-有限性对唯一性是必要的。\Cref{ex:2-3-2-5} 将在
$([0,1],\#)\times([0,1],m)$ 上给出两个测度：它们在每个矩形上都等于
$\#(A)m(B)$，在对角线上的质量却分别为 $0$ 与 $1$。因此存在性证明之前必须明确限制唯一性的范围。

\begin{lemma}[有限简单函数的质量和]\label{lem:product-simple-section-functional}
设 $\mu(X)<\infty$。若非负简单函数有不交表示
$g=\sum_{j=1}^rc_j\1_{A_j}$，定义
$I_\mu(g)=\sum_jc_j\mu(A_j)$。该数与不交表示无关。若
$0\le g_n\uparrow g$，所有函数均简单、有共同有限界，则
$I_\mu(g_n)\uparrow I_\mu(g)$。
\end{lemma}

\begin{proof}
两种表示用所有 $A_j\cap B_k$ 作共同有限分割；在每个非空交块上两边系数相同，故质量和相同。
设所有函数由 $M$ 控制。给定 $\varepsilon>0$，集合
$D_n=\{x:g(x)-g_n(x)<\varepsilon\}$ 递增至 $X$；这是因为每一点上的实数列递增至 $g(x)$。
测度从下连续和 $\mu(X)<\infty$ 给 $\mu(X\setminus D_n)\to0$。共同有限分割还给简单质量和的
有限线性：在 $g,g_n$ 的所有水平集交块上逐项相减，得到
$I_\mu(g-g_n)=I_\mu(g)-I_\mu(g_n)$。因而
\[
 0\le I_\mu(g)-I_\mu(g_n)=I_\mu(g-g_n)
 \le\varepsilon\mu(X)+M\mu(X\setminus D_n).
\]
先令 $n\to\infty$，再令 $\varepsilon\downarrow0$，得到所需收敛。
\end{proof}

\begin{theorem}[乘积测度存在唯一性]\label{theorem:2-3-2-2}
若 $\mu,\nu$ 都 $\sigma$-有限，则存在唯一乘积测度 $\mu\otimes\nu$。
\end{theorem}

\begin{proof}
先设两测度有限。可测矩形构成半环，其有限不交并构成环 $\mathcal R$。若
$R\in\mathcal R$，则 $R_x$ 是有限个可测集合之并，函数
$g_R(x)=\nu(R_x)$ 是非负简单函数。确切地，若
$R=\bigcup_{l=1}^r(A_l\times B_l)$，则由 $(A_l)_{l=1}^r$ 生成的至多 $2^r$ 个胞腔
\[
 C_J=\Bigl(\bigcap_{l\in J}A_l\Bigr)
      \cap\Bigl(\bigcap_{l\notin J}A_l^c\Bigr),
 \qquad J\subset\{1,\dots,r\},
\]
使 $R_x=\bigcup_{l\in J}B_l$ 在每个非空 $C_J$ 上固定。因此 $g_R$ 至多取 $2^r$ 个值。定义
\[
 \rho_0(R)=I_\mu(g_R).
 \label{eq:product-premeasure}
\]
共同矩形分割和上一引理说明定义与 $R$ 的表示无关，并且
$\rho_0(A\times B)=\mu(A)\nu(B)$。

若 $R_1,\dots,R_N$ 两两不交，则截面也两两不交，所以对每个 $x$，
\[
 g_{\bigsqcup_{k=1}^NR_k}(x)
 =\nu\!\left(\bigsqcup_{k=1}^N(R_k)_x\right)
 =\sum_{k=1}^N\nu((R_k)_x)
 =\sum_{k=1}^Ng_{R_k}(x).
\]
在这些简单函数水平集的共同有限分割上，$I_\mu$ 的有限线性性给
\[
 \rho_0\!\left(\bigsqcup_{k=1}^NR_k\right)
 =\sum_{k=1}^N\rho_0(R_k).
\]

若 $R=\bigsqcup_{k\ge1}R_k$ 且各集合都在环中，令
$S_N=\bigsqcup_{k\le N}R_k$。逐个 $x$ 有
$\nu((S_N)_x)\uparrow\nu(R_x)$；极限仍是有界简单函数，因为 $R$ 属于矩形环且 $\nu(Y)<\infty$。
有限简单单调引理给
\[
 \rho_0(R)=\lim_N\rho_0(S_N)=\sum_{k=1}^{\infty}\rho_0(R_k).
\]
所以 $\rho_0$ 是预测度。环版 \Caratheodory{} 扩张给
$\mathcal A\otimes\mathcal B$ 上的测度；有限测度唯一性给唯一。

一般情形取两两不交的可测层 $(X_i)_{i\ge1}$、$(Y_j)_{j\ge1}$，使
\[
 X=\bigsqcup_iX_i,\qquad Y=\bigsqcup_jY_j,\qquad
 \mu(X_i)<\infty,\quad \nu(Y_j)<\infty.
\]
在 $X_i\times Y_j$ 上对限制测度应用有限情形，所得乘积记为 $\rho_{i,j}$。对
$E\in\mathcal A\otimes\mathcal B$ 定义
\[
 \rho(E)=\sum_{i,j\ge1}\rho_{i,j}\bigl(E\cap(X_i\times Y_j)\bigr).
 \label{eq:sigma-finite-product-patching}
\]
若 $E_n$ 两两不交，则每个块内的限制仍两两不交；先用 $\rho_{i,j}$ 的可数可加性，再对非负三重
级数换序，得到
\[
 \rho\Bigl(\bigsqcup_nE_n\Bigr)
 =\sum_{i,j,n}\rho_{i,j}\bigl(E_n\cap(X_i\times Y_j)\bigr)
 =\sum_n\rho(E_n).
\]
所以 $\rho$ 是测度。对矩形 $A\times B$，
\[
\begin{aligned}
 \rho(A\times B)
 &=\sum_{i,j}\mu(A\cap X_i)\nu(B\cap Y_j)\\
 &=\left(\sum_i\mu(A\cap X_i)\right)
   \left(\sum_j\nu(B\cap Y_j)\right)
 =\mu(A)\nu(B),
\end{aligned}
\]
其中各项非负，$0\cdot\infty=0$ 的情形也由逐项为零直接得到。这证明全局存在性。

这里块限制的可测性可直接核对。固定 $i,j$，令 $\mathcal T_{i,j}$ 为所有满足
\[
 E\cap(X_i\times Y_j)
 \in(\mathcal A|_{X_i})\otimes(\mathcal B|_{Y_j})
\]
的 $E\subset X\times Y$。该类对补集和可数并封闭；并且
\[
 (A\times B)\cap(X_i\times Y_j)
 =(A\cap X_i)\times(B\cap Y_j),
\]
所以它包含所有可测矩形，进而包含 $\mathcal A\otimes\mathcal B$。因此
\eqref{eq:sigma-finite-product-patching} 的每一项都确实有定义。反过来，所有形如
$E\cap(X_i\times Y_j)$、$E\in\mathcal A\otimes\mathcal B$ 的集合构成块上的 $\sigma$-代数，
并包含每个 $(A\cap X_i)\times(B\cap Y_j)$。由两边的最小性，实际得到迹等式
\[
 (\mathcal A\otimes\mathcal B)|_{X_i\times Y_j}
 =(\mathcal A|_{X_i})\otimes(\mathcal B|_{Y_j}).
\]

最后令 $P_N=(\bigcup_{i\le N}X_i)\times(\bigcup_{j\le N}Y_j)$。这些矩形递增覆盖
$X\times Y$，且任一满足矩形公式的候选测度在 $P_N$ 上都有相同有限质量。可测矩形构成
$\pi$-系统并生成 $\mathcal A\otimes\mathcal B$，故
\cref{theorem:2-2-1-2} 的 $\sigma$-有限唯一性部分给出全局唯一性。
\end{proof}

\begin{example}[Lebesgue 乘积]\label{ex:2-3-2-1}
先在 Borel 定义域上陈述。若 $m_{d,B}$ 表示 $\R^d$ 上的 Borel Lebesgue 测度，则对半开盒
$A\subset\R^n$、$B\subset\R^k$，$m_{n,B}\otimes m_{k,B}$ 与 $m_{n+k,B}$ 在
$A\times B$ 上都取值 $m_{n,B}(A)m_{k,B}(B)$。有理端点开盒构成 Euclidean 拓扑的可数基，而
每个这样的开盒都是递增半开盒的可数并，故半开乘积盒生成 $\Borel(\R^{n+k})$。它们还给出共同
有限耗尽 $[-N,N)^n\times[-N,N)^k$。唯一性因此给出
\[
 m_{n,B}\otimes m_{k,B}=m_{n+k,B}
 \quad\text{在 }\Borel(\R^n)\otimes\Borel(\R^k)
 =\Borel(\R^{n+k})\text{ 上}.
\]
完成这一个 Borel 乘积测度便得到 $m_{n+k}$。这里没有把
$\mathcal L(\R^n)\otimes\mathcal L(\R^k)$ 与高维 Lebesgue 完成化擅自视为同一个
$\sigma$-代数；完成化与乘积定义域的进一步比较留到 Tonelli 理论建立之后。
\end{example}

\begin{example}[乘积平面中的对角线]\label{ex:2-3-2-2}
对角线 $\Delta=\{(x,x):x\in\R\}$ 闭，因而 Borel。固定 $M>0$，沿对角线用至多
$\lceil2M/\varepsilon\rceil+1$ 个边长 $2\varepsilon$ 的方盒覆盖
$\Delta\cap[-M,M]^2$。总面积至多
\[
 4\varepsilon^2\left(\frac{2M}{\varepsilon}+2\right)
 =8M\varepsilon+8\varepsilon^2.
\]
令 $\varepsilon\downarrow0$，该有界部分零测；再对 $M\in\N$ 取可数并，得到
$(m\otimes m)(\Delta)=0$。
\end{example}

\begin{example}[不可数乘积中的两种 $\sigma$-代数]\label{ex:2-3-2-3}
令 $I$ 不可数，$X=\{0,1\}^I$。柱 $\sigma$-代数中的每个集合只依赖某个可数坐标集：生成柱只
依赖有限坐标，而“依赖至多可数坐标”的集合类对补集和可数并封闭。全零点
$0^I$ 的单点集却在乘积拓扑中闭，因为
\[
 \{0^I\}=\bigcap_{i\in I}\{x:x_i=0\}.
\]
它不依赖任何可数坐标集：固定可数 $J\subset I$ 后，在 $I\setminus J$ 的一个坐标改成 $1$，
所得点与全零点在 $J$ 上相同。因此该闭集不在柱 $\sigma$-代数，乘积拓扑的 Borel
$\sigma$-代数严格更大。
\end{example}

\begin{example}[非 $\sigma$-有限时矩形规则不唯一]\label{ex:2-3-2-5}
令 $X=[0,1]$，第一因子取计数测度 $\#$，第二因子取 Lebesgue 测度 $m$。对 Borel
$E\subset X^2$，置
\[
 \lambda(E)=\sum_{x\in X}m(E_x),\qquad
 \eta(E)=m\{x:(x,x)\in E\},
\]
其中不可数非负和定义为所有有限部分和的上确界。对两两不交的 $E_n$，每个截面也不交；交换两个
非负求和次序得到
\[
\begin{aligned}
 \lambda\!\left(\bigsqcup_nE_n\right)
 &=\sum_{x\in X}m\!\left(\bigsqcup_n(E_n)_x\right)
  =\sum_{x\in X}\sum_nm((E_n)_x)\\
 &=\sum_n\sum_{x\in X}m((E_n)_x)
  =\sum_n\lambda(E_n).
\end{aligned}
\]
这里两个迭代和都等于所有有限集合
$F\subset X$ 与有限指标集 $G\subset\N$ 的矩形部分和
$\sum_{x\in F,n\in G}m((E_n)_x)$ 的上确界，故换序不需要预设不可数级数的条件收敛。
$\eta$ 是对角映射
$x\mapsto(x,x)$ 对 $m$ 的推前，因而也是测度。

矩形上
$\lambda(A\times B)=\#A\,m(B)$。而
$\eta(A\times B)=m(A\cap B)$。若 $m(B)=0$，则两式中的附加项为零，且约定
$\#A\,m(B)=0$；若 $m(B)>0$ 而 $A$ 至多可数，则 $m(A\cap B)=0$；若 $m(B)>0$ 且 $A$
不可数，则 $\#A\,m(B)=\infty$，加上不超过 $1$ 的 $m(A\cap B)$ 仍为无穷。因此
$\lambda$ 与 $\lambda+\eta$ 都满足同一矩形
规则。可是
\[
 \lambda(\Delta)=0,\qquad (\lambda+\eta)(\Delta)=1.
\]
计数测度不 $\sigma$-有限，正是唯一性失效的位置。
\end{example}

\begin{proposition}[矩形近似]\label{proposition:2-3-2-3}
若 $E\in\mathcal A\otimes\mathcal B$ 且
$(\mu\otimes\nu)(E)<\infty$，则对每个 $\varepsilon>0$，存在有限个可测矩形的并 $R$ 使
\[
 (\mu\otimes\nu)(E\mathbin\triangle R)<\varepsilon.
\]
\end{proposition}

\begin{proof}
先在有限乘积空间中证明。令 $\mathcal C$ 为可被矩形代数按对称差测度任意逼近的集合类。它含矩形
代数。若 $E_n\uparrow E$ 且 $E_n\in\mathcal C$，从下连续性使
$(\mu\otimes\nu)(E\setminus E_N)$ 对某个 $N$ 小于 $\varepsilon/2$，再逼近 $E_N$。递减情形由有限
总质量和补集化为递增情形：矩形代数含全空间且对补封闭，并且若 $R$ 逼近 $E$，则
\[
 E^c\mathbin\triangle R^c=E\mathbin\triangle R.
\]
故 $\mathcal C$ 对补集封闭，递减列取补后即化为递增列。因此 $\mathcal C$ 是含矩形代数的单调类，集合单调类定理给
$\mathcal C=\mathcal A\otimes\mathcal B$。

一般 $\sigma$-有限情形，先取共同有限矩形耗尽 $Q_n$，使
$(\mu\otimes\nu)(E\setminus Q_n)<\varepsilon/2$，再在 $Q_n$ 内使用有限情形。
\end{proof}

\begin{proposition}[拓扑 Borel 一致性]\label{proposition:2-3-2-4}
若 $X,Y$ 第二可数，则
\[
 \Borel(X\times Y)=\Borel(X)\otimes\Borel(Y).
\]
一般总有右侧包含于左侧。
\end{proposition}

\begin{proof}
坐标投影连续，所以可测矩形为 Borel 集，给出一般包含。第二可数时，取可数基
$(U_i),(V_j)$。每个乘积开集是它所含基矩形 $U_i\times V_j$ 的并；这样的基矩形只有可数多个，
所以开集属于右侧，给出反向包含。
\end{proof}

\begin{figure}[htbp]
\centering
\begin{tikzpicture}[scale=.9]
 \draw[book line] (-3,-1.8) rectangle (3,1.8);
 \fill[BookAccent!8] (-2.4,-1.3) rectangle (-.5,.8);
 \draw[BookAccent,semithick] (-2.4,-1.3) rectangle (-.5,.8);
 \node at (-1.45,-.25) {$A\times B$};
 \draw[BookWarning,thick] (1.2,-1.5) .. controls (2,-.8) and (.4,.2) .. (1.8,1.4);
 \draw[dashed] (1.25,-1.8)--(1.25,1.8) node[above] {$x$};
 \fill[BookWarning] (1.25,-.91) circle (1.5pt);
 \fill[BookWarning] (1.25,.42) circle (1.5pt);
 \node[right] at (1.35,-.2) {$E_x$};
 \draw[book accent arrow] (3.4,0)--(4.7,0) node[midway,above]{扩张};
 \node[book strong node] at (6.1,0) {$\mu\otimes\nu$\\全局测度};
\end{tikzpicture}
\caption{矩形规则先在环上形成预测度；截面把积空间集合重新变成参数化的单空间集合。}
\label{fig:product-sections-extension}
\end{figure}

\begin{remark}
对 $E\in\mathcal A\otimes\mathcal B$，每个截面可测；但函数
$x\mapsto\nu(E_x)$ 的可测性以及
$(\mu\otimes\nu)(E)$ 与其迭代积分的关系要到 Tonelli 定理才建立。不可数乘积中也必须区分柱
$\sigma$-代数与乘积拓扑的 Borel $\sigma$-代数。
\end{remark}

在概率论中，独立随机变量的联合分布由边缘分布的乘积刻画；把加法映射作用于乘积测度会产生卷积。
Kolmogorov 扩张定理则从相容的有限维分布构造随机过程的分布。两因子乘积的存在与矩形唯一性，是这些
后续构造以及 Tonelli 定理的共同起点。

\begin{exercise}
证明可测集合的截面可测，并说明反命题为什么失败。
\end{exercise}
\begin{exercise}
仅用矩形唯一性证明 $m_n\otimes m_k=m_{n+k}$。
\end{exercise}
\begin{exercise}
证明第二可数空间的乘积 Borel 等于 Borel 乘积。
\end{exercise}
\begin{exercise}
逐项核对非 $\sigma$-有限例中 $\lambda,\eta$ 的可数可加性及所有矩形情形。
\end{exercise}

乘积测度把多坐标观测组合起来；在实线上，次序还允许用半直线质量把整份测度压缩成一个单调函数。

\subsection{分布函数与 Lebesgue--Stieltjes 测度}\label{subsec:2-3-3}
\index{Lebesgue--Stieltjes 测度}\index{分布函数}

实线有一个其他空间没有的便利：集合可以按从左到右的次序累积。有限测度 $\mu$ 的半直线质量
$F(x)=\mu(( -\infty,x])$ 是单调函数。原子表现为跳跃，没有质量的开区间表现为平台。反过来，
若给定右连续非减函数，区间增量 $F(b)-F(a)$ 应当生成测度；关键在于证明这些有限增量对可数分解
连续。

历史上，Stieltjes 的这类增量构造与连分数及矩问题相伴发展；在这里，它的作用是把单调函数编码成
测度。第~4~章引入有界变差与 Riemann--Stieltjes 积分后，我们再比较“由增量生成测度”与“对函数
作 Stieltjes 积分”这两个方向。

\begin{definition}[分布函数]\label{def:2-3-3-1}
实线上有限 Borel 测度 $\mu$ 的分布函数为
\[
 F_\mu(x)=\mu(( -\infty,x]).
\]
若 $\mu$ 是概率，则 $F_\mu(-\infty)=0$、$F_\mu(+\infty)=1$。
\end{definition}

\begin{definition}[Stieltjes 增量]\label{def:2-3-3-2}
设 $F:\R\to\R$ 右连续且非减。在半开区间上置
\[
 \mu_F((a,b])=F(b)-F(a),\qquad a<b.
\]
无限端点只在相应单调极限存在且不产生 $\infty-\infty$ 时使用。下文假设
$F(-\infty)=\lim_{x\to-\infty}F(x)$ 有限。
\end{definition}

\begin{definition}[跳跃]\label{def:2-3-3-3}
非减函数在 $x$ 的跳跃为
\[
 \Delta F(x)=F(x)-F(x-),\qquad F(x-)=\lim_{t\uparrow x}F(t).
\]
\end{definition}

\begin{proposition}[单调函数的跳跃至多可数]\label{proposition:2-3-3-1}
实线上的单调函数只有至多可数多个跳跃点。
\end{proposition}

\begin{proof}
以非减函数为例。固定 $M,n\in\N$，在 $[-M,M]$ 中跳跃至少 $1/n$ 的不同点至多
\[
 n\bigl(F(M)-F(-M-)\bigr)+1
\]
个。事实上，若这样的点为 $-M\le x_1<\cdots<x_r\le M$，则对任意
$\varepsilon>0$，忽略最左边的一个点，并对 $j=1,\ldots,r-1$ 选取
$t_j\in(x_j,x_{j+1})$，使
\[
 x_1<t_1<x_2<t_2<x_3<\cdots<t_{r-1}<x_r
\]
并且 $F(x_{j+1})-F(t_j)>1/n-\varepsilon/r$；这里用的是
$F(t)\to F(x_{j+1}-)$ 当 $t\uparrow x_{j+1}$。这些实轴区间可取为两两不交，单调性于是给
\[
 \sum_{j=1}^{r-1}\bigl(F(x_{j+1})-F(t_j)\bigr)
 \le F(M)-F(-M),
\]
令 $\varepsilon\downarrow0$，得到
$(r-1)/n\le F(M)-F(-M)\le F(M)-F(-M-)$，从而得到所写的上界。
所有正跳跃点是这些有限集合对 $M,n$ 的可数并。非增情形对 $-F$ 使用同一论证。
\end{proof}

\begin{lemma}[Stieltjes 环预测度从上连续]\label{lem:stieltjes-continuity-at-empty}
在有界半开区间及其有限不交并组成的环上，以区间增量定义的有限可加函数 $\mu_0$ 满足：若
$E_n\downarrow\varnothing$，则 $\mu_0(E_n)\downarrow0$。因此 $\mu_0$ 是预测度。
\end{lemma}

\begin{proof}
若 $E=\bigsqcup_{i=1}^r(a_i,b_i]$，先置
\[
 \mu_0(E)=\sum_{i=1}^r\bigl(F(b_i)-F(a_i)\bigr).
\]
若 $E$ 还有另一有限不交区间表示，收集两种表示的全部端点并细分；在每个分量上恒等式
\[
 F(b)-F(a)=\sum_{l=1}^N\bigl(F(c_l)-F(c_{l-1})\bigr),
 \qquad a=c_0<\cdots<c_N=b,
\]
说明两种总和都等于同一组小区间增量之和。因此 $\mu_0$ 良定义，且同一共同细分立即给出有限
可加性。设
$E_n\downarrow\varnothing$，并反设
$\inf_n\mu_0(E_n)=\eta>0$。把每个 $E_n$ 写成有限个不交区间 $(a,b]$。由 $F$ 在每个左端点
右连续，若第 $n$ 轮有 $r_n$ 个非空分量 $(a_{n,l},b_{n,l}]$，可在每个分量中选
$c_{n,l}\in(a_{n,l},b_{n,l})$，使
\[
 F(c_{n,l})-F(a_{n,l})<\frac{\eta2^{-n-2}}{r_n}.
\]
故各左端薄条 $(a_{n,l},c_{n,l}]$ 的增量总和小于 $\eta2^{-n-2}$。令 $K_n$ 为相应闭区间
$[c_{n,l},b_{n,l}]$ 的有限并，令 $H_n$ 为这些薄条的有限并。
于是 $K_n\subset E_n$ 是紧集，$H_n$ 属于该环，且 $H_n$ 覆盖 $E_n\setminus K_n$，并有
$\mu_0(H_n)<\eta2^{-n-2}$。这里没有给未必属于半开区间环的差集
$E_n\setminus K_n$ 直接赋值。

令 $L_n=\bigcap_{j=1}^nK_j$。若某个 $L_n$ 为空，则每个 $x\in E_n$ 至少被前 $n$ 轮中的一个
薄条删去，所以 $E_n\subset\bigcup_{j\le n}H_j$。有限次可加与次可加给
\[
 \mu_0(E_n)\le\sum_{j=1}^n\mu_0(H_j)<\eta/2,
\]
其中第一个不等式只用环上的非负有限可加性。确切地，令
\[
 H'_1=H_1,
 \qquad H'_j=H_j\setminus\bigcup_{i<j}H_i\quad(2\le j\le n).
\]
则 $(E_n\cap H'_j)_{j\le n}$ 两两不交并覆盖 $E_n$，所以
\[
 \mu_0(E_n)=\sum_{j=1}^n\mu_0(E_n\cap H'_j)
 \le\sum_{j=1}^n\mu_0(H'_j)
 \le\sum_{j=1}^n\mu_0(H_j).
\]
与定义 $\mu_0(E_n)\ge\eta$ 矛盾。因此每个 $L_n$ 非空。它们是递减非空紧集，并包含于有界紧集
$L_1$，有限交性质给
$\bigcap_nL_n\ne\varnothing$；但该交包含于 $\bigcap_nE_n=\varnothing$，再次矛盾。

最后说明预测度。若环中 $A=\bigsqcup_{k\ge1}A_k$，令
$R_N=A\setminus\bigsqcup_{k\le N}A_k$。则 $R_N\downarrow\varnothing$，有限可加性给
$\mu_0(A)=\sum_{k\le N}\mu_0(A_k)+\mu_0(R_N)$。令 $N\to\infty$ 即得可数可加。
\end{proof}

\begin{theorem}[Lebesgue--Stieltjes 构造]\label{theorem:2-3-3-2}
每个右连续非减函数 $F:\R\to\R$，若 $F(-\infty)$ 有限，则唯一确定局部有限 Borel 测度
$\mu_F$，满足
\[
 \mu_F((a,b])=F(b)-F(a),\qquad
 \mu_F(( -\infty,b])=F(b)-F(-\infty).
\]
给 $F$ 加常数不改变 $\mu_F$。
\end{theorem}

\begin{proof}
上一引理给出半开有界区间环上的预测度。环版 \Caratheodory{} 扩张产生 Borel 测度；每个有界
区间的质量有限，所以它局部有限；每个紧集包含于有界区间，故紧集有限。实线第二可数且局部紧，
\cref{theorem:2-2-2-2} 还说明该测度为 Radon。半开区间在耗尽 $(-n,n]$ 上形成共同有限的 $\pi$-系统，
\cref{theorem:2-2-1-2} 给出唯一性。对
$(a,b]$ 令 $a\to-\infty$，测度从下连续给第二个公式。所有增量只含函数值之差，故加常数不改变
测度。
\end{proof}

\begin{corollary}[紧区间常值延拓版]\label{cor:stieltjes-compact-interval}
设 $G:[a,b]\to\R$ 右连续且非减，给定 $\alpha\ge0$。定义
\[
 F(x)=
 \begin{cases}
 G(a)-\alpha,&x<a,\\
 G(x),&a\le x\le b,\\
 G(b),&x>b.
 \end{cases}
\]
则 $F$ 生成支撑包含于 $[a,b]$ 的有限测度，并且
\[
 \mu_F(\{a\})=\alpha,
 \qquad \mu_F((s,t])=G(t)-G(s)\quad(a\le s<t\le b).
\]
取 $\alpha=0$ 就是两端常值延拓且左端无原子；若要保留指定的左端原子，必须把左侧常值降低
$\alpha$。这个端点约定在以后把区间上的 BV 函数变成 Stieltjes 测度时保持不变。
\end{corollary}

\begin{proof}
分段函数右连续且非减，并在两端外常值，故上一构造适用。由测度从上连续和构造定理的区间公式，
\[
 \mu_F(\{a\})
 =\lim_{n\to\infty}\mu_F((a-1/n,a])
 =\lim_{n\to\infty}\bigl(F(a)-F(a-1/n)\bigr)=\alpha.
\]
内部半开区间公式就是 Stieltjes 增量。补集
$\R\setminus[a,b]$ 的两部分可逐项计算：
\[
 \mu_F(( -\infty,a))
 =\lim_{n\to\infty}\bigl(F(a-1/n)-F(-\infty)\bigr)=0,
\]
并且由从下连续性
\[
 \mu_F((b,\infty))
 =\lim_{n\to\infty}\mu_F((b,n])
 =\lim_{n\to\infty}\bigl(F(n)-F(b)\bigr)=0.
\]
所以测度集中于 $[a,b]$。
\end{proof}

\begin{proposition}[跳跃就是原子]\label{prop:stieltjes-jump-atom}
对 Stieltjes 测度，
\[
 \mu_F(\{x\})=\Delta F(x).
\]
\end{proposition}

\begin{proof}
集合 $(x-1/n,x]$ 递减至 $\{x\}$，且第一项质量有限。测度从上连续给
\[
 \mu_F(\{x\})=\lim_n\bigl(F(x)-F(x-1/n)\bigr)=F(x)-F(x-).
\]
\end{proof}

\begin{example}[纯原子分布]\label{ex:2-3-3-1}
若 $\mu=\sum_kp_k\delta_{x_k}$、$p_k>0$ 且 $\sum_kp_k<\infty$，则
\[
 F_\mu(x)=\sum_{k:x_k\le x}p_k.
\]
点 $x$ 的跳跃为所有满足 $x_k=x$ 的权重之和。若把有理数枚举为 $(q_k)$ 并取
$p_k=2^{-k}$，跳跃点在每个区间稠密，$F$ 没有非平凡平台，支撑为整条实线。纯原子因此不等于
支撑离散。
\end{example}

\begin{example}[均匀分布]\label{ex:2-3-3-2}
$[0,1]$ 上均匀概率的分布函数是
\[
 F(x)=\begin{cases}0,&x<0,\\x,&0\le x\le1,\\1,&x>1.
 \end{cases}
\]
它连续，且除两个端点外导数为 $\1_{(0,1)}$；导数不是在整条实线上恒为 $1$。
\end{example}

\begin{example}[Cantor 分布]\label{ex:2-3-3-3}
标准 Cantor 函数的逐层构造及连续性、单调性证明明确推迟到
\cref{ex:2-4-3-1}。使用那里将证明的性质：$C$ 连续非减，从 $0$ 增长到 $1$，并在 Cantor 集补集的每个三分开区间上
恒定。它的 Stieltjes 测度没有原子，因为 $C$ 没有跳跃；补集是可数个恒定区间之并，每个区间
质量为零，所以全部质量集中在零 Lebesgue 测度的 Cantor 集上。$C$ 在该补集上局部恒定，因而
$C'=0$ 在补集每点成立；Cantor 集零测，所以导数几乎处处为零，而总增量仍为一。导数遗漏了这份
奇异连续质量；相应分解在后章建立。
\end{example}

\begin{theorem}[测度—分布函数对应]\label{theorem:2-3-3-3}
有限 Borel 测度 $\mu$ 与满足右连续、非减、有界且 $F(-\infty)=0$ 的函数一一对应；总质量为
$F(+\infty)$。
\end{theorem}

\begin{proof}
给定有限 $\mu$，半直线嵌套给 $F_\mu$ 非减。若 $x_j\downarrow x$，则
$(-\infty,x_j]\downarrow(-\infty,x]$；有限测度从上连续给右连续。令
$x\to-\infty$、$x\to+\infty$，分别用从上连续于空集和从下连续于全空间，得到端点值
$0,\mu(\R)$。区间差给
\[
 \mu((a,b])=F_\mu(b)-F_\mu(a).
\]
反过来，对规范化 $F$ 使用 Stieltjes 构造。所得测度满足
$\mu_F(( -\infty,b])=F(b)-F(-\infty)=F(b)$，所以其分布函数就是 $F$。若两个有限 Borel 测度
有同一分布函数，则它们在所有 $(a,b]$ 上相等，而且由从下连续性
\[
 \mu(\R)=\lim_{n\to\infty}F(n)=\nu(\R).
\]
半开区间生成 $\Borel(\R)$，有限测度唯一性因而给出 $\mu=\nu$。所以两次操作互逆。
\end{proof}

\begin{proposition}[连续点上的分布函数收敛]\label{proposition:2-3-3-4}
设 $\mu_n,\mu$ 是 $\R$ 上的有限正 Borel 测度，分布函数分别为 $F_n,F$。若
$F_n(x)\to F(x)$ 对 $F$ 的每个连续点成立，则
\[
 \int f\,\dd\mu_n\longrightarrow\int f\,\dd\mu
 \qquad(f\in C_c(\R)).
\]
若还有 $\mu_n(\R)\to\mu(\R)$，则同一结论对每个 $f\in C_b(\R)$ 成立。因此，总质量条件
对紧支撑测试可以删除，但这类测试可能看不见逃向无穷的质量。
\end{proposition}

证明见\cref{corollary:7-3-2-finite-measures}。前面的 \(\delta_n\) 说明，只有紧支撑测试函数时，总质量可能逸向
无穷。

\begin{proposition}[平台读出支撑]\label{proposition:2-3-3-5}
对有限 Borel 测度 $\mu$，
\[
 \R\setminus\supp\mu
 =\bigcup\{(a,b):F_\mu\text{ 在 }(a,b)\text{ 上恒定}\}.
\]
等价地，$x\in\supp\mu$ 当且仅当 $F_\mu$ 在 $x$ 的每个开邻域内都有正增量。
\end{proposition}

\begin{proof}
若 $F$ 在 $(a,b)$ 恒定，则每个有理端点子区间 $(c,d]\subset(a,b)$ 的质量
$F(d)-F(c)$ 为零。这些子区间可数地覆盖 $(a,b)$，故整段零测，其中每点都不在支撑中。
反过来，若 $x\notin\supp\mu$，存在开区间 $(a,b)\ni x$ 且 $\mu((a,b))=0$。对
$a<c<d<b$，有 $F(d)-F(c)=\mu((c,d])=0$，所以 $F$ 在该区间恒定。
\end{proof}

\begin{figure}[htbp]
\centering
\begin{tikzpicture}[x=1.1cm,y=2.2cm]
 \draw[->] (-.3,0)--(6.3,0) node[right]{$x$};
 \draw[->] (0,-.08)--(0,1.25) node[above]{$F(x)$};
 \draw[BookAccent,thick] (0,.08)--(1.3,.08);
 \draw[BookAccent,thick] (1.3,.08)--(2.8,.55);
 \draw[BookAccent,thick] (2.8,.55)--(3.6,.55);
 \draw[BookAccent,thick] (3.6,.78)--(4.5,.78);
 \draw[BookAccent,thick] (4.5,.78).. controls (5,.9) and (5.4,1.08)..(6,1.12);
 \draw[BookWarning,dashed] (3.6,.55)--(3.6,.78);
 \node[below] at (.65,0){零质量平台};
 \node[above] at (2,.36){连续增长};
 \node[right] at (3.6,.67){原子};
 \node[above] at (5.2,.95){弥散增长};
\end{tikzpicture}
\caption{分布函数中的平台、连续增长和跳跃分别编码零质量区、弥散质量与原子。}
\label{fig:stieltjes-platforms-jumps}
\end{figure}

\begin{remark}
全书固定右连续函数和 $(a,b]$ 增量；换成左连续函数时半开方向也必须同时更换。$F$ 与 $F+C$
生成同一测度，只有规定 $F(-\infty)=0$ 才与有限测度一一对应。几乎处处导数不能恢复 Cantor
分布的奇异连续部分。
\end{remark}

分布函数还产生分位数：对给定概率水平寻找最小的 \(x\)，使 \(F(x)\) 达到该水平。它也把矩问题变成
“若干积分是否足以识别测度”的问题；Cantor 例说明，几乎处处导数不能恢复全部测度。分位数、矩与
弱收敛将在积分理论建立后继续讨论。

\begin{exercise}
证明单调函数的跳跃点至多可数，并给出跳跃点稠密的纯原子分布。
\end{exercise}
\begin{exercise}
计算有限离散分布和一个离散—均匀混合分布的分布函数。
\end{exercise}
\begin{exercise}
从给定右连续阶梯函数构造测度，并核对所有半开区间增量。
\end{exercise}
\begin{exercise}
用从上连续性证明 $\mu_F(\{x\})=F(x)-F(x-)$。
\end{exercise}

至此，测度已经能够从有限集合数据被构造、由生成类识别，并沿映射和乘积传输。下一节回到
Euclidean 几何，建立第 3 章换元公式所需的局部线性模型和零集控制。
